\setlength{\footskip}{8mm}


\chapter{DISCUSSION}
The development and evaluation of the ALIGN-VQA architecture expose fundamental truths about how artificial intelligence must approach longitudinal medical imaging. For years, the prevailing methodology in Difference Visual Question Answering (Diff-VQA) has relied on the implicit assumption that deep neural networks can naturally discern temporal changes through raw feature subtraction or concatenation. The findings of this thesis categorically dismantle that assumption. This chapter synthesizes the empirical results, examines the unexpected behaviors of the model, justifies the critical architectural decisions that drove the success, and candidly addresses the boundaries of the current approach.

\section{Revisiting Research Questions}
The central thesis of this work was that longitudinal VQA requires explicit spatial alignment, targeted noise filtration, and question-guided reasoning. Across the three research questions, the empirical data not only validated these claims but revealed unexpected dynamics regarding how vision-language models process temporal medical data.

\subsection{Cross Image Difference Attention}
The author fully expected that semantically aligning the reference image to the current image would improve performance. However, the magnitude of the baseline’s failure without difference attention module was highly unexpected. When processing un-aligned scans using direct subtraction, the model’s linguistic output almost entirely collapsed, dropping to a BERT-F1, \cite{bert-score}, of 0.818 and CIDEr, \cite{cider}, of 1.383. This catastrophic drop revealed a critical insight. Without explicit spatial alignment, generative language models do not simply make minor diagnostic errors. They lose the ability to ground their language in visual reality, resorting to fragmented, incoherent guesses. The module proved that spatial correspondence is not merely a performance booster although it is a fundamental prerequisite for coherent longitudinal reasoning.

\subsection{Learnable Residual Gating}
The author hypothesized that explicitly filtering the raw subtraction features would remove static background noise and improve diagnostic accuracy. However, what the author did not expect was the emergent semantic sophistication of the gate itself. The author initially assumed the gating mechanism would function primarily as a rigid spatial mask, simply zeroing out unchanged bones or static medical hardware. Qualitative analysis of the gated features revealed that the model learned to handle highly complex, dynamic non-pathological variations. For instance, it actively suppressed visual artifacts caused by differing levels of patient inspiration, deep of a breath the patient took between the two scans. As a profound realization, the network didn't just learn to mask unchanged pixels although it learned a deep clinical heuristic for successfully distinguishing between a shift in lung volume  and a genuine change in lung fluid/opacity.

\subsection{Question Guided Tokens and Targeted Reasoning}
The author anticipated that pruning tokens would make the model more computationally efficient although the author expected a slight trade-off in diagnostic accuracy due to lost information. The results completely subverted this expectation. The sparse Top-16 token configuration actually outperformed the dense Top-49 configuration CIDEr 2.427 vs. 2.226. This is a profound finding. It suggests a paradox, less is more, in medical VQA: dense visual features act as a distraction. Because the vast majority of a chest X-ray remains unchanged between visits, feeding the language decoder 49 tokens overwhelms its cross-attention layers with static, irrelevant anatomy. By forcing the Question Guided Tokenizer (QDT) to aggressively prune 70 percent of the image and pass forward only 16 highly relevant tokens, it inadvertently improved the decoder’s focus. However, the author also found the lower bound, which pushes the sparsity to Top-12 tokens resulted in a performance drop, indicating the threshold where critical contextual margins around a pathology are destroyed.

\section{Comparing with Past Work}
ALIGN-VQA successfully achieved the primary claim, which is a lightweight, highly accurate architecture capable of interpreting disease progression while filtering out non-pathological noise. ALIGN-VQA currently outperforms state-of-the-art models like \cite{real} and \cite{plural}, achieving peak structural and semantic scores. However, ALIGN-VQA did not yet achieve flawless reasoning on sub-pixel or highly abstract relational questions involving complex spatial reasoning across multiple intersecting pathologies. The ultimate goal is to develop a fully autonomous, zero-shot longitudinal reasoning engine that matches the contextual awareness of an experienced radiologist, a system capable of tracking disease trajectories over not just two, but dozens of sequential visits. We are currently one major step closer to this goal; by solving the spatial misalignment and noise bottlenecks, the author has laid the necessary groundwork. The remaining distance requires shifting from static two-image comparison to dynamic, multi-hop temporal graph reasoning.

\section{Experimental Decisions}
The success of ALIGN-VQA is heavily predicated on specific, deliberate experimental and architectural design choices. Rather than relying on standard transformer pipelines, the author engineered the workflow to mirror the cognitive process of a radiologist.

\subsection{Architectural Sequencing}
A major architectural decision was the strict sequential ordering of the modules. In early ideation, one might assume that pruning tokens first and then aligning them would save compute. The author deliberately chose to align the dense features first, gate the dense residuals, and only then apply the QDT to prune the tokens.

This decision was critical to the results. Spatial alignment requires full global context as the network must see the clavicles, the diaphragm, and the rib cage to calculate the semantic transformation matrices. If the model had pruned the tokens prior to alignment based on the text question, the CIDA module would be starved of the global anatomical landmarks required to align the images. Similarly, the Learnable Residual Gate needs the surrounding healthy tissue to establish a baseline for normal noise before it can filter it. By placing the QDT at the very end of the visual encoder pipeline, the author ensured that the tokens selected for the language decoder were already perfectly aligned and scrubbed of acquisition artifacts. This design choice directly enabled the high CIDEr scores observed in the ablation studies.

\subsection{The Selection of NLG Metrics for Clinical Evaluation}
A second crucial experimental decision was the heavy reliance on CIDEr and ROUGE-L, \cite{rouge}, alongside traditional BLEU, \cite{bleu}, metrics. In standard computer vision, accuracy or simple BLEU scores are often deemed sufficient. The author explicitly designed the evaluation around CIDEr because medical reporting is a long-tail linguistic distribution. In a clinical setting, a generic phrase like "The lungs are clear" appears thousands of times in the dataset, whereas a critical phrase like "mild pneumothorax has worsened" appears rarely. BLEU treats all words equally. If a model defaults to safe, high-frequency phrasing, it can artificially inflate its BLEU score while missing the actual disease. The author chose CIDEr because its Term Frequency-Inverse Document Frequency (TF-IDF) weighting heavily penalizes models that guess generic templates and explicitly rewards the model for generating rare, highly specific diagnostic keywords. By utilizing CIDEr as the primary metric for diagnostic accuracy, the author forced the gating mechanism to genuinely learn the difference between background noise and true, rare pathological evolution, rather than allowing the model to cheat the evaluation through language priors.

\section{Limitation and Future Work}
While ALIGN-VQA demonstrates state-of-the-art performance in longitudinal reasoning, it is constrained by a fundamental structural limitation regarding how it perceives multi-dimensional space.

\subsection{Limitation}
The most concrete limitation affecting the validity of this architecture is its absolute reliance on 2-Dimensional feature projections. While the CIDA module is highly effective at correcting 2D spatial misalignments, such as planar translations, vertical slouching, or varying inspiration volumes, it cannot semantically reconstruct 3D out-of-plane rotations. 

In clinical reality, a patient might be slightly rotated laterally relative to the X-ray detector in their current visit compared to their reference visit. Because a chest X-ray is a 2D projection of a 3D volume, an out-of-plane rotation fundamentally alters the overlapping densities of the ribs, heart, and lung parenchyma. As the CIDA module operates on 2D semantic maps, it will attempt to forcefully align these features, which may warp or distort the underlying pathological representation. Consequently, the Learnable Residual Gate might misinterpret the sudden appearance of a rotated rib shadow as a "new opacity," or conversely, gate out a true lesion because the rotational alignment obscured it. This 3D-to-2D projection ambiguity represents a hard ceiling on the accuracy of any purely 2D longitudinal subtraction model, including ALIGN-VQA.

\subsection{Future Work}
Future iterations of ALIGN-VQA should abandon 2D constraints and simultaneously process paired posteroanterior and lateral views. By constructing a multi-view temporal graph, the architecture could triangulate the Z-axis depth of a pathology. If the CIDA module is extended to cross-reference alignment matrices between posteroanterior and lateral views, the model could mathematically differentiate between an out-of-plan patient rotation, which affects both views systematically, and a genuine volumetric tumor growth, effectively nullifying the current core limitation.

Currently, the Learnable Residual Gate filters visual noise before the model processes the clinical question. While this works well for general noise reduction, what constitutes "noise" is often dependent on the clinical context. For example, surgical hardware is "noise" if the question asks about pneumonia, but it is the "signal" if the question asks about a pacemaker lead. Future work should implement a Question-Conditioned Gate, where the text embedding of the clinical query is injected directly into the gating bottleneck. This would allow the gate to dynamically redefine its filtration parameters on a per-query basis, resulting in even higher efficiency and diagnostic precision.

\begin{comment}

\chapter{DISCUSSION}
This is perhaps the most important parts, because it holds the insight!  It shows how smart and wise you are as a researcher.

\textbf{Write at least two pages.}

So many things you could have discussed.  This is guideline:

\section{Revisiting Research Questions}
\begin{itemize}
    \item Did you find what you expect?  Discuss in depth anything you did not expect.
    \item Did you achieve what you claim?  and what you did not yet achieve? what is your ultimate goal, and how far are you from it? 
\end{itemize}
    
\section{Comparing with Past Work}
\begin{itemize}
    \item How did your work compare with related work? consistent, anything surprising? 
\end{itemize}
    
\section{Experimental Decisions}
\begin{itemize}
    \item Why did you choose certain experimental/approaches design decisions, and how did it affect your results?
    \item You can specifically talk about certain aspects of your work 
\end{itemize}
    
\section{Design Implications}
Design implications or guidelines.   This is more tailored for HCI study.

\section{Limitations and Future Work}
\begin{itemize}
    \item Limitations: for limitations, please only describe the core limitation that you think might affect the validity of your result. Limitation such as “I conduct this study in the US but I have not tested in the UK” is not a limitation to be discussed as it is almost a common limitation in any study.  Discuss CONCRETE, REAL limitation.
    \item Future work, you only have to discuss one or two really concrete one.
\end{itemize}

\end{comment}